\section{N-gram-mallit}

N-gram-malli on klassinen tilastollinen kielimalli, jossa seuraavan sanan tai symbolin ennustaminen perustuu sitä edeltävään  kontekstiin. Mallin taustalla on $(n-1)$:n kertaluvun Markovin oletus, jonka mukaan nykyinen tokeni riippuu vain edeltävistä $(n-1)$ tokenista eikä koko aiemmasta historiallisesta kontekstista. Tällöin voidaan approksimoida ehdollinen todennäköisyys

\begin{equation*}
  P(w_t | w_1, w_2, \ldots, w_{t-1}) \approx P(w_t | w_{t-n+1}, \ldots, w_{t-1}),
\end{equation*}
missä $w_t$ on ajanhetken $t$ tokeni ja $(w_{t-n+1}, \ldots, w_{t-1})$ on sitä edeltävät $(n-1)$ tokenia.

Olkoon $V=\{v_1, v_2, \ldots, v_{|V|}\}$ sanasto, jossa |V| on sanaston koko ja $W=(v_{i_1}, v_{i_2}, \ldots, v_{i_{|W|}})\\=(w_1, w_2, \ldots, w_{|W|})$ havaintoaineisto. N-gram-malli voidaan toteuttaa muodostamalla lukumatriisi $C\in\Re^{|V|^{n-1} \times |V|}$, jossa $C_{c,j}$ edustaa kuinka monta kertaa kontekstia $c$ seuraa tokeni $v_j$ havaintoaineistossa. Formaalisti voidaan kirjoittaa
\begin{align*}
  C_{c,j} = |\{t\in\N | (w_{t-n+1}, \dots, w_{t-1}) = c \text{ ja } w_t = v_j \}|.
\end{align*}

Kun lukumatriisin $C$ rivit normalisoidaan jakamalla jokainen alkio rivinsä summalla, saadaan frekvenssimatriisi $F\in\Re^{|V|^{n-1} \times |V|}$, jossa
\begin{align*}
  F_{c,j} = \frac{C_{c,j}}{\sum_{k=1}^{|V|} C_{c,k}},
\end{align*} edustaa ehdollista todennäköisyyttä $P(w_t = v_j | (w_{t-n+1}, \dots, w_{t-1}) = c)$. On helppo osoittaa, että nämä ehdolliset todennäköisyydet vastaavat suurimman uskottavuuden estimaatteja.

N-gram-mallin oletuksen mukaisesti koko havaintoaineiston todennäköisyys voidaan kirjoittaa ehdollisten todennäköisyyksien tulona
\begin{equation*}
  P(W) \approx \prod_{t=1}^{|W|} P(w_t \mid w_{t-n+1}, \ldots, w_{t-1}).
\end{equation*}
Kiinteän kontekstin $c$ tapauksessa havaintojen lukumäärät
\[
  \{C_{c,1}, C_{c,2}, \ldots, C_{c,|V|}\}
\]
muodostavat multinomijakauman otoksen, jossa parametreina ovat todennäköisyydet
\[
  \theta_c = (\theta_{c,1}, \theta_{c,2}, \ldots, \theta_{c,|V|}),
  \qquad
  \sum_{j=1}^{|V|}\theta_{c,j} = 1.
\]
Tällöin kyseisen kontekstin osalta uskottavuusfunktio voidaan kirjoittaa muodossa
\begin{equation}\label{eq:n-gram_target_function}
  \mathcal{L}(\theta_c)
  = \prod_{j=1}^{|V|} \theta_{c,j}^{C_{c,j}}.
\end{equation}
Tämän funktion maksimoiminen on ekvivalenttia sen logaritmin maksimoimiselle, joten saadaan Lagrangen funktio
\begin{align*}
  \mathcal{L}_{\text{Lag}}(\theta_c, \lambda)
  &= \sum_{j=1}^{|V|} C_{c,j} \log \theta_{c,j} + \lambda\left(1 - \sum_{j=1}^{|V|} \theta_{c,j}\right).
\end{align*}
Derivoimalla saadaan
\begin{align*}
  &\frac{\partial \mathcal{L}_{\text{Lag}}}{\partial \theta_{c,j}}
  = \frac{C_{c,j}}{\theta_{c,j}} - \lambda = 0,
  && j=1,2,\ldots,|V| \\
  &\implies \theta_{c,j} = \frac{C_{c,j}}{\lambda}, && j=1,2,\ldots,|V| \\
  &\implies \lambda = \sum_{j=1}^{|V|} C_{c,j} \\
  &\implies \theta_{c,j} = \frac{C_{c,j}}{\sum_{k=1}^{|V|} C_{c,k}}, && j=1,2,\ldots,|V|.
\end{align*}
eli täsmälleen rivikohtaisesti normalisoitu frekvenssimatriisi.

Tekstin generointi n-gram-mallilla tapahtuu iteratiivisesti hyödyntämällä frekvenssimatriisia $F$. Ensin valitaan alkukonteksti, esimerkiksi erikoismerkeillä aloitettu $(n-1)$-tokenin pituinen alkujakso. Tämän jälkeen mallin tämänhetkistä kontekstia vastaava rivi $F_{c,\cdot}$ tulkitaan todennäköisyysjakaumaksi, josta seuraava tokeni valitaan satunnaisotannalla. Kun tokeni $w_t$ on generoitu, konteksti päivitetään siirtämällä ikkuna yhden askeleen eteenpäin siten, että uusi konteksti muodostuu edeltävistä $(n-1)$ generoidusta tokenista. Prosessia jatketaan, kunnes generoidaan päätemerkki tai saavutetaan haluttu tekstin pituus.

Formaalisti voidaan sanoa, että jos nykyinen konteksti on $c = (w_{t-n+1}, \ldots, w_{t-1})$, niin seuraava tokeni valitaan jakaumasta
\begin{align*}
  w_t \sim F_{c,\cdot}
\end{align*}  ja uusi konteksti on $c' = (w_{t-n+2}, \ldots, w_t)$.
Näin muodostunut teksti on siis peräkkäisten ehdollisten otosten ketju, jossa kunkin tokenin valinta riippuu vain edeltävästä paikallisesta kontekstista.

Vaikka lähestymistapa on käsitteellisesti yksinkertainen ja tulkittava, siihen liittyy merkittäviä laskennallisia haasteita erityisesti kontekstikoon \(n\) kasvaessa. Tämän seurauksena lukumatriisista $C\in\Re^{|V|^{n-1} \times |V|}$ tulee nopeasti erittäin suuri ja käytännöllisissä sovelluksissa hyvin harva, koska vain pieni osa kaikista mahdollisista konteksteista esiintyy tyypillisessä havaintoaineistossa. Tämä lisää sekä muistin- että laskentatehon tarvetta.

Kuvattujen rajoitusten vuoksi lukumatriisiin perustuva menetelmä ei skaalaudu hyvin suuriin \(n\)-arvoihin. Erityisesti harva havaintotiheys ja eksponentiaalinen tilavaatimuksen kasvu rajoittavat mallin käytännöllisyyttä. Tästä syystä on perusteltua tarkastella vaihtoehtoisia menetelmiä, joissa mallin parametrit opitaan optimoimalla minimoitavaa häviöfunktiota.

\subsection{N-gram-malli optimointiongelmana}

Sen sijaan, että määriteltäisiin mallin parametrit suoraan laskemalla havaintofrekvenssejä, voidaan lähestyä ongelmaa maksimointiperiaatteen kautta. Tällöin tavoitteena on löytää mallin parametrit $\theta$, jotka maksimoivat havaintoaineiston uskottavuuden $P(W;\theta)$.

Koneoppimisessa optimointiongelma esitetään yleensä minimointimuodossa. Tämän vuoksi on tavanomaista määritellä häviöfunktioksi negatiivinen log-uskottavuus. Jotta mallin häviötä voidaan vertailla eri havaintoaineistojen välillä, jaetaan häviöfunktio havaintoaineiston koolla eli
\begin{equation*}
  \mathcal{J}(\theta; W)
  = -\frac{1}{|W|}\sum_{t=1}^{|W|}
  \log P_\theta(w_t \mid w_{t-n+1}, \ldots, w_{t-1})
\end{equation*}

Tämä suure on keskimääräinen negatiivinen log-todennäköisyysarvo, jonka minimointi on täsmälleen ekvivalenttia uskottavuuden maksimoinnille. Toisin sanoen
\begin{equation*}
  \arg\max_{\theta} P(W;\theta)
  =
  \arg\min_{\theta} \mathcal{J}(W;\theta).
\end{equation*}

Tämä näkökulma on keskeinen, sillä se muuntaa probabilistisen mallinnuksen optimointitehtäväksi, jota voidaan ratkaista yleisillä numeerisen optimoinnin menetelmillä.

Neuroverkkopohjaisessa $n$-gram-mallissa tavoitteena on oppia funktio
\begin{equation*}
  f_{\theta} : \mathcal{C} \to [0,1]^{|V|},
\end{equation*}
jossa $\mathcal{C}$ on kontekstien avaruus ja $\theta$ mallin opittavat parametrit. Malli saa syötteenä kontekstin $c=(w_{t-n+1},\ldots,w_{t-1})$ ja tuottaa jakauman seuraavalle tokenille:
\begin{equation*}
  f_{\theta}(c) = p_{\theta}(\cdot \mid c).
\end{equation*}
Tällöin halutun tokenin todennäköisyys on
\begin{equation*}
  p_{\theta}(w_t \mid c),
\end{equation*}
ja koko aineiston häviöfunktio voidaan kirjoittaa muodossa
\begin{equation*}
  \mathcal{J}(\theta)
  = - \frac{1}{|W|}\sum_{t=1}^{|W|} \log p_{\theta}(w_t \mid w_{t-n+1}, \ldots, w_{t-1}).
\end{equation*}

\subsection{Neuroverkot optimointiongelmana}

Edellä kuvattu n-gram-malli perustuu suoraan havaittuihin frekvensseihin tai niiden
maksimointiin perustuvaan estimaattiin. Tällöin mallin parametrit määräytyvät
käsin määritellyistä tilastollisista säännöistä, eikä mallilla ole kykyä oppia
kontekstien välisiä monimutkaisempia riippuvuuksia. Neuroverkkopohjaisessa
lähestymistavassa taas pyritään oppimaan suoraan funktio, joka kuvaa annetun
kontekstin ja sitä seuraavan tokenin välistä yhteyttä.

Yleisesti neuroverkko voidaan esittää peräkkäisten parametrisoitujen funktioiden
kompositiona. Olkoon \(x \in \mathbb{R}^d\) syötevektori ja \(y \in \mathbb{R}^m\)
mallin tuottama ulostulo. Tällöin yksi piilokerros voidaan kirjoittaa muodossa

$$
h = \phi(Wx + b),
$$

missä \(W \in \mathbb{R}^{k \times d}\) on painomatriisi, \(b \in \mathbb{R}^k\) on
siirtovektori ja \(\phi\) on epälineaarinen aktivaatiofunktio, kuten ReLU,
tanh tai sigmoid. Monikerroksisessa verkossa tämä prosessi toistetaan useita
kertoja, jolloin saadaan

$$
h^{(1)} = \phi^{(1)}(W^{(1)}x + b^{(1)}), \quad
h^{(2)} = \phi^{(2)}(W^{(2)}h^{(1)} + b^{(2)}), \ \ldots
$$

ja lopulta verkon tuottama lopputulos saadaan viimeisestä kerroksesta.

Kielimallinnuksessa neuroverkon tehtävänä on arvioida todennäköisyysjakauma
seuraavalle tokenille annetun kontekstin ehdolla. Jos konteksti on
\(c = (w_{t-n+1}, \ldots, w_{t-1})\), pyritään oppimaan malli

$$
f_\theta : C \rightarrow [0,1]^{|V|},
$$

missä \(f_\theta(c)\) on kontekstin \(c\) perusteella muodostettu todennäköisyysjakauma
sanaston kaikille tokeneille. Käytännössä tämä toteutetaan usein siten, että
konteksti muunnetaan numeeriseksi syötevektoriksi \(x_c\), jonka perusteella
verkko tuottaa ensin niin sanotut logit-arvot \(z \in \mathbb{R}^{|V|}\). Näistä
muodostetaan todennäköisyysjakauma softmax-funktiolla

$$
p_\theta(w_t = v_j \mid c) =
\frac{\exp(z_j)}{\sum_{k=1}^{|V|} \exp(z_k)}.
$$

Tällöin mallin oppimistavoite voidaan esittää negatiivisena log-uskottavuuden
keskiarvona, eli ristientropiahäviönä:

$$
J(\theta) = -\frac{1}{|W|} \sum_{t=1}^{|W|}
\log p_\theta(w_t \mid w_{t-n+1}, \ldots, w_{t-1}).
$$

Tämä on sama tavoitefunktio kuin n-gram-mallin tapauksessa, mutta nyt
todennäköisyydet eivät määräydy suoraan frekvenssien perusteella vaan ne opitaan
parametrisella funktiolla \(f_\theta\). Tällöin optimointiongelmana on löytää
parametrit \(\theta\), jotka minimoivat häviöfunktion

$$
\theta^\ast = \arg\min_\theta J(\theta).
$$

Neuroverkoissa tämän minimoinnin suorittaminen perustuu gradienttipohjaisiin
optimointimenetelmiin, kuten stokastiseen gradienttilaskuun
(stochastic gradient descent, SGD) ja sen muunnelmiin, esimerkiksi Adam-
optimointiin. Parametreja päivitetään laskemalla häviöfunktion gradientti
suhteessa verkon painoihin ja siirtoihin sekä siirtämällä parametreja tämän
gradientin vastakkaiseen suuntaan. Tämä prosessi tunnetaan takaisinsyöttönä
(backpropagation), ja sen avulla verkko oppii asteittain parantamaan seuraavan
tokenin ennustamista.

Keskeinen etu neuroverkkopohjaisessa mallissa on, että samaan aikaan kun se
mallintaa kontekstin ja seuraavan tokenin välistä yhteyttä, se oppii myös
jatkuvia esityksiä sanoille ja konteksteille. Toisin kuin lukumatriisipohjainen
n-gram-malli, neuroverkko ei ole sidottu vain havaittuihin konteksteihin, vaan
se voi yleistää osittain myös harvinaisempiin tai kokonaan näkemättömiin
tilanteisiin oppimiensa painojen ja esitysten avulla. Tästä syystä neuroverkkoja
voidaan pitää n-gram-ajattelun yleistyksenä, jossa tilastollinen taulukko korvataan
opittavalla funktiolla.
